\section{Revisi\'on de literatura}

La literatura reciente sobre sue\~no con EEG monocanal muestra dos tendencias relevantes para este proyecto. La primera es que la mayor parte del rendimiento competitivo en staging proviene de \textbf{representaciones espectrales y de complejidad} bien dise\~nadas, incluso cuando el clasificador final no es profundo \cite{gao2023psdrf,nakamura2017complexity,ghimatgar2019hmm,li2022cascadedsvm}. La segunda es que, para apnea, el problema no es solo la clasificaci\'on binaria sino la explicaci\'on fisiol\'ogica de qu\'e bandas del EEG sostienen la decisi\'on \cite{barnes2022apneaeegcnn}. Esa segunda observaci\'on es especialmente relevante aqu\'i, porque la detecci\'on de apnea con un solo canal EEG es m\'as fr\'agil que con PSG multimodal.

Desde el punto de vista metodol\'ogico, varios trabajos convergen en cuatro ideas. Primero, la ventana de 30 s sigue siendo la unidad can\'onica por su compatibilidad con el etiquetado de sue\~no y con la sincronizaci\'on de eventos respiratorios. Segundo, las \emph{features} m\'as persistentes son potencia por bandas, potencias relativas, relaciones espectrales y medidas de complejidad o entrop\'ia \cite{gao2023psdrf,nakamura2017complexity}. Tercero, los modelos secuenciales o con contexto temporal suelen mejorar N1 y las transiciones, pero ese beneficio no siempre se sostiene cuando el protocolo es estrictamente \emph{subject-wise} \cite{mousavi2019sleepeegnet,ghimatgar2019hmm}. Cuarto, la comparaci\'on entre papers rara vez es totalmente justa, porque cambia el canal EEG, la definici\'on de etiquetas, el preprocesamiento, el balance de clases y, en muchos casos, el protocolo de partici\'on.

\begin{table*}[t]
\centering
\caption{Trabajos de literatura que guiaron el preprocesamiento y la selecci\'on de \emph{features}.}
\label{tab:literature_preprocessing}
\resizebox{\textwidth}{!}{%
\begin{tabular}{lllll}
\toprule
Referencia & Canal / se\~nal & Segmentaci\'on & Rasgos metodol\'ogicos principales & Relevancia para este trabajo \\
\midrule
\cite{mousavi2019sleepeegnet} & EEG monocanal (Sleep-EDF) & 30 s + contexto secuencial & Representaci\'on tiempo-frecuencia y contexto entre \'epocas & Justifica usar contexto y explica por qu\'e N1 mejora con dependencia temporal \\
\cite{barnes2022apneaeegcnn} & EEG monocanal & Ventanas respiratorias / EEG & Importancia de bandas delta y beta para apnea & Sustenta la selecci\'on de \feature{eeg\_bandpower\_*} y potencias relativas para apnea \\
\cite{gao2023psdrf} & EEG monocanal/dual & 30 s & PSD de ondas caracter\'isticas + Random Forest & Baseline m\'as cercano a nuestra l\'inea cl\'asica de staging \\
\cite{nakamura2017complexity} & EEG monocanal & 30 s con vecinos temporales & Entrop\'ias multiescala + SEF + SVM & Motiv\'o incluir entrop\'ias y SEF en la l\'inea multitarget \\
\cite{ghimatgar2019hmm} & EEG monocanal & 30 s & HMM sobre rasgos EEG & Sirve como referencia secuencial no implementada \\
\cite{li2022cascadedsvm} & EEG monocanal & 30 s & SVM en cascada y normalizaci\'on cuidadosa & Relevante para interpretar el papel de SVM en nuestro baseline \\
\bottomrule
\end{tabular}%
}
\end{table*}


La Tabla~\ref{tab:literature_preprocessing} resume los rasgos metodol\'ogicos que fueron m\'as influyentes para nuestro pipeline: uso de EEG monocanal, \'epocas fijas, potencia espectral, entrop\'ias y an\'alisis de comparabilidad. A partir de esa revisi\'on, este proyecto adopt\'o como n\'ucleo de \emph{features} variables de potencia absoluta/relativa, medidas de complejidad, estad\'isticos de forma y descriptores de frecuencia de borde espectral.

\begin{table*}[t]
\centering
\caption{Baselines publicados y grado de comparabilidad con el protocolo del proyecto.}
\label{tab:literature_baselines}
\resizebox{\textwidth}{!}{%
\begin{tabular}{lllllll}
\toprule
Referencia & Dataset & Canal & Tarea & Modelo & Comparable & Nota \\
\midrule
\cite{mousavi2019sleepeegnet} & Sleep-EDF & EEG monocanal & Sleep staging & Seq2Seq deep & No & Usa deep learning; comparable solo por dataset y tarea \\
\cite{barnes2022apneaeegcnn} & Apnea cohort & EEG monocanal & Apnea & CNN explicable & Parcial & Comparable por tarea y canal, no por familia de modelo \\
\cite{gao2023psdrf} & Sleep-EDF & Fpz-Cz / Pz-Oz & Sleep staging & PSD + Random Forest & Si & Muy cercano a nuestro baseline de staging \\
\cite{nakamura2017complexity} & Sleep EEG & EEG & Sleep staging & Complejidad/HMM-style & Parcial & Comparable por monocanal y features, no por pipeline final \\
\cite{ghimatgar2019hmm} & Sleep EEG & EEG monocanal & Sleep staging & HMM & Parcial & Sirve como contraste secuencial, no como linea principal \\
\cite{li2022cascadedsvm} & Sleep EEG & EEG monocanal & Sleep staging & SVM en cascada & Si & Muy comparable con la familia SVM del proyecto \\
\cite{wang2024svmxgboost} & Sleep EEG & EEG & Sleep staging & SVM + XGBoost & Si & Comparable por familias clasicas y objetivo \\
\cite{satapathy2024multimodality} & Multimodal sleep data & Multimodal & Sleep staging & ML multimodal & No & Sirve como referencia externa, pero no es justo por usar mas senales \\
\bottomrule
\end{tabular}%
}
\end{table*}


La Tabla~\ref{tab:literature_baselines} clasifica los baselines publicados por grado de comparabilidad con nuestro protocolo. Los trabajos m\'as cercanos al baseline de staging son el enfoque PSD + Random Forest sobre Sleep-EDF \cite{gao2023psdrf}, el enfoque de complejidad + SVM \cite{nakamura2017complexity} y la propuesta de SVM en cascada \cite{li2022cascadedsvm}. Para apnea, el referente m\'as cercano por tarea y restricci\'on de canal es Barnes \emph{et al.} \cite{barnes2022apneaeegcnn}. En contraste, trabajos multimodales o enteramente profundos son \'utiles como contexto, pero no como comparaci\'on estricta con la l\'inea principal de este informe.
